\section{Recursive Traversals: Preorder, Inorder, Postorder}
\label{sec:trees-rec-traversals}

\problemstatement{Given the root of a binary tree, return its preorder,
inorder, and postorder traversals. Preorder visits \emph{node, left,
right}; inorder visits \emph{left, node, right}; postorder visits
\emph{left, right, node}.}

\begin{patternbox}
\emph{``Visit every node in a fixed relative order''} is the definition of
a depth-first traversal. The only thing separating the three is
\textbf{where the current node is printed relative to its two recursive
calls} -- before both (pre), between them (in), or after both (post). If
you can write one, you can write all three by moving a single line.
\end{patternbox}

\subsection*{Understanding the problem carefully}

The recursive structure of a binary tree hands you the traversal almost
for free. To process a whole tree you process its root and its two
subtrees; each subtree is processed the same way. The word
``pre/in/post'' refers only to \emph{when the root itself is recorded}
relative to recursing into the left and right subtrees.

\begin{ideabox}[One Template, Three Positions]
Every depth-first traversal is the same three actions -- \textsc{record
node}, \textsc{recurse left}, \textsc{recurse right} -- performed in a
different order:
\begin{itemize}
  \item \textbf{Preorder:} record, left, right.
  \item \textbf{Inorder:} left, record, right.
  \item \textbf{Postorder:} left, right, record.
\end{itemize}
The recursion itself is byte-for-byte identical; only the placement of the
``record'' step moves.
\end{ideabox}

\subsection*{Algorithm}

\begin{enumerate}
  \item \textsc{Inorder}(\textit{node}, \textit{out}):
  \begin{enumerate}
    \item If \textit{node} is null, return.
    \item \textsc{Inorder}(\textit{node}.left, \textit{out}).
    \item Append \textit{node}.val to \textit{out}.
    \item \textsc{Inorder}(\textit{node}.right, \textit{out}).
  \end{enumerate}
  \item Preorder and postorder are the same three lines with the append
        moved to the top or the bottom respectively.
\end{enumerate}

\subsection*{Dry run}

\dryrun{Tree: root $1$, right child $2$, and $2$'s left child $3$.}

\begin{itemize}
  \item \textbf{Preorder}: record $1$; left of $1$ is null; go right to
        $2$; record $2$; left of $2$ is $3$; record $3$ (its children
        null); right of $2$ null. Result: $1,2,3$.
  \item \textbf{Inorder}: left of $1$ null; record $1$; go right to $2$;
        left of $2$ is $3$, record $3$; record $2$; right null. Result:
        $1,3,2$.
  \item \textbf{Postorder}: left of $1$ null; right subtree of $1$: left
        of $2$ is $3$, record $3$; right of $2$ null; record $2$; record
        $1$. Result: $3,2,1$.
\end{itemize}

\subsection*{C++ implementation}

\begin{lstlisting}
#include <bits/stdc++.h>
using namespace std;

struct TreeNode {
    int val;
    TreeNode* left;
    TreeNode* right;
    TreeNode(int v) : val(v), left(nullptr), right(nullptr) {}
};

void preorder(TreeNode* node, vector<int>& out) {
    if (node == nullptr) return;
    out.push_back(node->val);   // record
    preorder(node->left, out);
    preorder(node->right, out);
}

void inorder(TreeNode* node, vector<int>& out) {
    if (node == nullptr) return;
    inorder(node->left, out);
    out.push_back(node->val);   // record
    inorder(node->right, out);
}

void postorder(TreeNode* node, vector<int>& out) {
    if (node == nullptr) return;
    postorder(node->left, out);
    postorder(node->right, out);
    out.push_back(node->val);   // record
}

int main() {
    TreeNode* root = new TreeNode(1);
    root->right = new TreeNode(2);
    root->right->left = new TreeNode(3);

    vector<int> pre, in, post;
    preorder(root, pre);
    inorder(root, in);
    postorder(root, post);

    auto printv = [](const string& name, vector<int>& v) {
        cout << name << ": ";
        for (int x : v) cout << x << " ";
        cout << "\n";
    };
    printv("preorder ", pre);
    printv("inorder  ", in);
    printv("postorder", post);
    return 0;
}
\end{lstlisting}

\begin{complexitybox}
\textbf{Time Complexity.} $O(n)$ -- each node is visited exactly once.
\textbf{Space Complexity.} $O(h)$ for the recursion stack, where $h$ is
the tree's height ($O(n)$ in the worst case of a skewed tree, $O(\log n)$
when balanced), plus $O(n)$ for the output array.
\end{complexitybox}

\begin{takeawaybox}
The three depth-first traversals are one algorithm with the ``record''
step in three different positions. Every later depth-first problem in this
chapter is a traversal with extra work grafted onto that record step, so
internalise this skeleton before moving on.
\end{takeawaybox}
